\begin{problem}[Gluing with an empty overlap]\label{prob:vi20-p16}
Prove gluing two schemes along the empty opens gives their disjoint union as a scheme.
\end{problem}

\ifdefined\FullProblemDossiers
\paragraph{Why this problem matters.}
This checks an extreme boundary case of P1.

\paragraph{Useful ideas before starting.}
No points or sections are forced to agree across charts.

\paragraph{Guided hints.}
\begin{enumerate}
\item No points or sections are forced to agree across charts.
\end{enumerate}

\begin{solution}
There are no cross-chart identifications, so the quotient space is the disjoint union.  Likewise, the compatibility condition on sections is vacuous across the two charts.  Sections therefore form the expected product on disjoint opens, which is precisely the disjoint-union sheaf.
\end{solution}

\paragraph{What happened mathematically.}
Disjoint unions are gluing with no overlap.

\paragraph{Extensions.}
Generalize to arbitrary families.

\paragraph{Provenance.}
\textbf{New canonical synthesis; no numbered legacy Problem is claimed.}
\else
\paragraph{Provenance.} \textbf{New canonical synthesis; no numbered legacy Problem is claimed.}
\fi
